\chapter{Experiments and Results}
\label{ch:exresults}

% ---------------------------------------------------------------------------
\section{Overview and Chapter Structure}
\label{sec:story-overview}
% ---------------------------------------------------------------------------

This chapter reports the experiments and results of the thesis. Each experiment
was motivated by the findings of the previous one, forming a chain from a
generative baseline to a corrected discriminative model and an assessment of
where it stands relative to the published KVP10k benchmark.

The experiments proceed as follows. We begin with a generative reference
baseline (Mistral-7B) and then evaluate a layout-aware discriminative approach
(LayoutLMv3). We then investigate two linking granularities, identify a
structural imbalance problem, and through diagnostic experiments identify and
correct three compounding data-pipeline bugs. Under the released KVP10k
benchmark, the final corrected model reaches regular-KVP F1 of 0.334
(text-only) and 0.309 (text+location), below the corresponding published values
of 0.659 and 0.611; diagnostic analyses localise the remaining gap to relation
modeling.

\subsection*{Evaluation Setup}

All official benchmark comparisons use the same evaluation set:
\textbf{581 unique test pages}
from the KVP10k test split (55\% of 1{,}051 unique pages; limited by PDF
download failures and scanned pages without text layers). The evaluation
follows IBM's released KVP10k benchmark implementation. Matching is greedy and
one-to-one, requires equal KVP types, and uses NED $<0.2$ for text. The paper
describes the location criterion as IoU $>0.3$, whereas the released code uses
IoU $\geq0.3$; benchmark-comparison results follow the executable code and
record this discrepancy. Precision and recall are calculated per page and
macro-averaged over pages with non-empty ground truth after category filtering;
F1 is the harmonic mean of the two macro averages. Results are reported for
regular KVPs, unkeyed values, unvalued keys, and their aggregate, under
text-only, location-only, and combined matching. Pooled regular-link counts
used during V1--V4 diagnosis are identified separately and are not compared
directly with the published macro scores; those diagnostics use explicitly
stated subsets and are not full-test benchmark results. Preprocessing and
input-text differences are described in Section~\ref{sec:data-preprocessing}.

% ---------------------------------------------------------------------------
\section{Experiment 1: Generative Baseline (Mistral-7B QLoRA)}
\label{sec:exp1}
% ---------------------------------------------------------------------------

\subsection*{Setup}

The starting point of the experimental work is a generative re-implementation
of the approach used in the KVP10k paper; it is not a strict reproduction.
Mistral-7B-Instruct-v0.2 is fine-tuned with 4-bit QLoRA (NF4 quantisation,
LoRA rank $r=4$, $\alpha=4$) on LMDX-formatted document prompts, where each
word in the document is represented as \texttt{word x1|y1|x2|y2} in
normalised coordinates. The model generates KVP predictions as structured text
output and is trained for 8 epochs with learning rate $5 \times 10^{-4}$.

The exact prompt template is shown in Figure~\ref{fig:stage3-prompt}. It follows
the LMDX paradigm~\cite{lmdx} used by the KVP10k baseline --- a document
representation with per-word coordinates followed by a task instruction
requesting a Python list-of-lists output. Because Naparstek et al.\ do not
publish their prompt, we reproduce its format and output contract rather than a
verbatim string. The reconstruction differs from the published system in
several respects. Word text and coordinates come from PyMuPDF's native text
layer rather than Tesseract OCR; pages without a usable text layer are
excluded; the resulting train and test subsets differ; the unpublished prompt
is reconstructed; LoRA rank is limited to $r=4$; and the paper does not report
enough training detail to reproduce its setup exactly. Numerical proximity
must therefore be read as an indicative comparison, not evidence of strict
reproduction.

\begin{figure}[ht]
\centering
\begin{minipage}{0.9\textwidth}
\small
\begin{verbatim}
<Document>
{lmdx_text}
</Document>
<Task>
From the document, extract the text keys and values.
Please provide the response in the form of a Python list of lists.
Each inner list should contain exactly two strings: the first string
is the key and the second string is the value.
Each key and value should be followed by its bounding box in the
format: left|top|right|bottom.
</Task>
### Response:
\end{verbatim}
\end{minipage}
\caption{Stage~3 LMDX prompt template. \texttt{\{lmdx\_text\}} is replaced by the
per-word text-and-coordinate representation of the page.}
\label{fig:stage3-prompt}
\end{figure}

\subsection*{Result 1: Competitive text-only F1, but unsupported generation}
\label{sec:res1}

\begin{table}[ht]
  \centering
  \caption{Experiment~1: pair-level macro F1 under the released KVP10k benchmark ($n=581$ available test pages). Published values are from Naparstek et al.\ (2024), Table~1.}
  \label{tab:exp1-overall}
  \begin{tabular}{@{}llccc@{}}
    \toprule
    \textbf{Category} & \textbf{System} & \textbf{Text} & \textbf{Location} & \textbf{Text+location} \\
    \midrule
    Regular  & KVP10k paper & 0.659 & 0.650 & 0.611 \\
             & Mistral (ours) & 0.598 & 0.547 & 0.521 \\
    \midrule
    Unkeyed  & KVP10k paper & 0.601 & 0.653 & 0.584 \\
             & Mistral (ours) & 0.000 & 0.000 & 0.000 \\
    \midrule
    Unvalued & KVP10k paper & 0.601 & 0.618 & 0.588 \\
             & Mistral (ours) & 0.744 & 0.699 & 0.672 \\
    \midrule
    All      & KVP10k paper & 0.643 & 0.661 & 0.612 \\
             & Mistral (ours) & 0.513 & 0.478 & 0.455 \\
    \bottomrule
  \end{tabular}
\end{table}

For regular KVPs, the reconstructed Mistral baseline reaches text-only
F1 $=0.598$, compared with $0.659$ in the paper, and combined F1 $=0.521$,
compared with $0.611$. Its unvalued-key score is strong, but the reconstructed
output parser emits no unkeyed-value predictions, giving zero for that category
and reducing the All score to $0.513$ text-only and $0.455$ combined. These
results are substantially different from the earlier pooled entity-level
calculation and replace it as the benchmark comparison. Given the data,
text-source, prompt, adapter-rank, and training-detail differences above, they
show that the reconstruction is competitive on regular pairs but do not
constitute a reproduction of the paper's baseline.

\paragraph{Per-cluster breakdown.}
\label{sec:stage3-per-cluster}
The fitted Stage 2 model assigns each coordinate-bearing test page using
the training-fitted scaler and K-means model. Table~\ref{tab:exp1-cluster}
recomputes the official regular-pair metrics from saved predictions and ground
truth; it does not rerun or retrain Mistral.

\begin{table}[ht]
  \centering
  \caption{Experiment~1 regular-KVP macro F1 by page-level
  annotation-geometry cluster. ``Scored'' counts pages with non-empty regular
  ground truth, following the released benchmark.}
  \label{tab:exp1-cluster}
  \begin{tabular}{@{}lrrrr@{}}
    \toprule
    \textbf{Stratum} & \textbf{Assigned} & \textbf{Scored} & \textbf{Text F1} & \textbf{Text+location F1} \\
    \midrule
    Cluster 0 (higher count, broad spread) & 319 & 304 & 0.651 & 0.569 \\
    Cluster 1 (lower count, compact)       & 213 & 179 & 0.498 & 0.429 \\
    Geometry unavailable                    &  49 &   0 & ---   & ---   \\
    \bottomrule
  \end{tabular}
\end{table}

Regular-pair performance is higher in Cluster~0 than Cluster~1, but the groups
describe annotation-box geometry rather than OCR density or document genre. A
separate pooled error audit shows unsupported-prediction rates of 37.1\% in
Cluster~0 and 53.0\% in Cluster~1. The 49 geometry-unavailable prepared pages
contain no regular ground truth and are therefore excluded from the official
per-cluster F1 calculation. These pages must not be interpreted as a third
performance regime. The overall benchmark values in Table~\ref{tab:exp1-overall}
are unchanged because only the grouping was recomputed.

\textbf{Motivation for Experiment 2:} A model that reads spatial positions
directly -- rather than generating approximate coordinates -- should be able
to predict only tokens that exist in the OCR input, preventing the generation
of unsupported text.

% ---------------------------------------------------------------------------
\section{Experiment 2: Layout-Aware Entity Classification (Stage 4a)}
\label{sec:exp2}
% ---------------------------------------------------------------------------

\subsection*{Setup}

We introduce LayoutLMv3-base (125M parameters) as the document encoder.
The reported runs use extracted text and bounding boxes, with no
document-specific page image content; the image input is a blank placeholder.
In Stage~4a, only the entity classification head is trained (linker disabled):
for each token, the model predicts whether it is a Key, Value, or Other token.
This phase establishes strong entity representations before the more complex
linking task is introduced.

\subsection*{Result 2: Strong Entity Classification}
\label{sec:res2}

\begin{table}[ht]
  \centering
  \caption{Experiment~2: Stage~4a entity classification results (best epoch).}
  \label{tab:exp2-results}
  \begin{tabular}{@{}lccc@{}}
    \toprule
    \textbf{Model} & \textbf{Precision} & \textbf{Recall} & \textbf{Entity F1} \\
    \midrule
    LayoutLMv3-base (entity only) & --- & --- & \textbf{0.847} \\
    \bottomrule
  \end{tabular}
\end{table}

Entity F1 = 0.847 confirms that LayoutLMv3 effectively classifies tokens as
keys, values, or other. Unlike Mistral, the discriminative model can only
predict labels for tokens that actually appear in the document -- there is
no mechanism for generating text that does not exist in the OCR input. It can,
however, still produce false positives by assigning an incorrect label to an
existing token. The resulting entity representations provide the foundation
for the linking stage.

\textbf{Motivation for Experiment 3:} With good entity classification in place,
the next challenge is linking: pairing each predicted key span with its
corresponding value span.

% ---------------------------------------------------------------------------
\section{Experiment 3: Token-Level Linking (V1)}
\label{sec:exp3}
% ---------------------------------------------------------------------------

\subsection*{Setup}

We activate the biaffine linker, trained jointly with the entity classifier.
In V1, linking operates at the \textbf{token level}: each key token is scored
against each value token, producing a dense adjacency matrix of size
$N_{key} \times N_{value}$. On a typical page with $\sim$100 key tokens and
$\sim$100 value tokens, this produces $\sim$10{,}000 candidate pairs with
only $\sim$20 true positives ($\sim$450:1 neg:pos ratio). A
\texttt{pos\_weight} of up to 50 is applied to the binary cross-entropy loss
to counteract the imbalance. A $\lambda$ sweep over $\{0.5, 1.0, 2.0\}$
controls the relative weight of the linking objective.

\subsection*{Result 3: Token-level class imbalance prevents linking}
\label{sec:res3}

\begin{table}[ht]
  \centering
  \caption{Experiment~3: V1 token-level linking results ($\lambda$ sweep).}
  \label{tab:exp3-results}
  \begin{tabular}{@{}lccc@{}}
    \toprule
    \textbf{$\lambda$} & \textbf{Entity F1} & \textbf{Link F1} & \textbf{TP / Pred / GT} \\
    \midrule
    0.5 & 0.849 & 0.007 & --- \\
    1.0 & 0.845 & 0.008 & 15 / 200 / 3{,}776 \\
    2.0 & 0.848 & 0.006 & --- \\
    \bottomrule
  \end{tabular}
\end{table}

Link F1 remains near zero across all $\lambda$ values. The model predicts only
200 links against 3{,}776 ground-truth links,\footnote{The 3{,}776 ground-truth links are counted over the full 581-page evaluation test set. This is a different data slice from the 200-row exploratory annotation subset (755 links) analysed in Chapter~\ref{ch:dataset} and from the per-cluster counts in Table~\ref{tab:exp8-cluster}; these are not conflicting counts of the same set.} of which only 15 are correct.
The $\lambda$ sweep confirms this is not a hyperparameter issue: the failure
is \textbf{structural}. With a $\sim$450:1 neg:pos ratio, even aggressive
class-weighting cannot force the linker to make positive predictions reliably.

\textbf{Motivation for Experiment 4:} The problem is the \emph{granularity}
of the linking candidates. Linking should not happen at the token level but
at the span level -- grouping contiguous key/value tokens into spans before
scoring would reduce the candidate space from $\sim$10{,}000 to $\sim$225
entries and the imbalance from 450:1 to 5:1.

% ---------------------------------------------------------------------------
\section{Experiment 4: Span-Level Linking (V2)}
\label{sec:exp4}
% ---------------------------------------------------------------------------

\subsection*{Setup}

In V2, contiguous same-label tokens are grouped into spans, and the linker
scores key spans against value spans. This reformulation cuts the number of
candidate pairs per page from roughly 10{,}000 to about 225 and, correspondingly,
reduces the negative-to-positive ratio from about 450:1 to about 5:1.
The original V2 implementation warm-started from the best Stage~4a checkpoint,
but span grouping used \texttt{argmax} predicted entity labels rather than
ground-truth labels; the effect of this choice is examined in
Experiments~5--7.

\subsection*{Result 4: Imbalance resolved, but link F1 remains near zero}
\label{sec:res4}

\begin{table}[ht]
  \centering
  \caption{Experiment~4: V2 span-level linking results.}
  \label{tab:exp4-results}
  \begin{tabular}{@{}lccc@{}}
    \toprule
    \textbf{Model} & \textbf{Entity F1} & \textbf{Link F1} & \textbf{TP / Pred / GT} \\
    \midrule
    V2 (span-level) & 0.847 & 0.007 & 37 / 498 / 3{,}776 \\
    \bottomrule
  \end{tabular}
\end{table}

Despite the favourable class balance, link F1 is still near zero. The
granularity fix was necessary but not sufficient. This unexpected result
triggered a systematic diagnostic investigation to find the root cause.

% ---------------------------------------------------------------------------
\section{Experiments 5--7: Diagnostic Analysis (D1, D2, D3)}
\label{sec:exp567}
% ---------------------------------------------------------------------------

\subsection*{Setup}

Three targeted diagnostic experiments were run to identify why V2 failed
despite correct class balance:

\paragraph{D1 --- Link score distribution.}
Raw sigmoid scores extracted from the best V2 checkpoint on 200 test
documents. Goal: determine whether the linker produces any high-confidence
predictions or is uniformly near-zero.

\paragraph{D2 --- Entity span analysis.}
Comparing ground-truth vs.\ predicted entity token counts before and after
bounding-box validity filtering in the span grouping function, per document.

\paragraph{D3 --- Link label audit.}
Tracing the full link label generation pipeline from ground-truth KVPs through
word-level labels, token-level alignment, and span-level collapse.

\subsection*{Result 5--7: Three Compounding Bugs Identified}
\label{sec:res567}

\paragraph{Bug 1 --- Missing Teacher Forcing.}
The V2 linker received \texttt{argmax} predicted entity labels (noisy) instead
of ground-truth labels for span grouping. Span boundaries were corrupted by
classification errors, misaligning the linker's input with the link labels.
\emph{Fix: pass ground-truth entity labels to the linker during training.}
Outcome: entity F1 = 0.847, link F1 = 0.007 --- unchanged, indicating deeper issues.

\paragraph{Bug 2 --- Bounding Box Filter (D2).}
The \texttt{group\_contiguous\_spans} function applied the strict predicate
$(x_1 < x_2) \wedge (y_1 < y_2)$, but LayoutLMv3's normalisation to the
$[0, 1000]$ coordinate grid produces many zero-width tokens ($x_1 = x_2$).
The filter rejected \textbf{49\% of key tokens} and \textbf{46\% of value
tokens}, cutting entity spans roughly in half. Only padding tokens
($[0,0,0,0]$) should be excluded.
\emph{Fix: filter only $[0,0,0,0]$ bounding boxes (CLS, SEP, PAD tokens).}

\paragraph{Bug 3 --- Link Label Cross-Product Explosion (D3).}
The label function \texttt{\_generate\_labels} created a link label for
\emph{every} key-word index $\times$ \emph{every} value-word index within a
KVP. For a typical page with 8 KVPs, this produced $\sim$1{,}049 positive
entries instead of 8. After span-level collapse, nearly every key--value pair
was marked positive, providing no discrimination for the linker to learn from
(equivalent to a uniform-positive supervision signal).
\emph{Fix: generate exactly one positive link label per ground-truth KVP pair.}

\begin{table}[ht]
  \centering
  \caption{Link label pipeline before and after V4 corrections (20 test docs, 179 GT KVPs).}
  \label{tab:diagnostic-verification}
  \begin{tabular}{@{}lcc@{}}
    \toprule
    \textbf{Metric} & \textbf{Before (V2)} & \textbf{After (V4)} \\
    \midrule
    Words per document & $\sim$35 & $\sim$200 \\
    Token-level link positives (20 docs) & $\sim$20{,}000 & 1{,}013 \\
    Word-level link recovery & 100\%* & 79\% (141/179) \\
    Span-level link recovery & 100\%* & 70\% (125/179) \\
    \bottomrule
  \end{tabular}
  \smallskip
  \footnotesize{*Before V4, ``100\% recovery'' is misleading: the cross-product produced all possible positive links, not just correct ones. Effective supervision was a near-uniform positive matrix -- equivalent to no supervision.}
\end{table}

Table~\ref{tab:diagnostic-verification} is a diagnostic verification summary,
not an ablation, and therefore cannot attribute effects to any single
correction. In V4, all three corrections are applied jointly. The reported V4
results therefore demonstrate pipeline repair as a whole rather than each
correction's independent numerical contribution; isolating individual effects
requires dedicated ablations outside the scope of this thesis.

% ---------------------------------------------------------------------------
\section{Experiment 8: Corrected Model (V4)}
\label{sec:exp8}
% ---------------------------------------------------------------------------

\subsection*{Setup}

V4 applies all three corrections jointly, is warm-started from the prior
\texttt{stage4b\_canary\_B} checkpoint, and is trained with:
\begin{itemize}[nosep]
  \item Teacher forcing enabled
  \item Corrected bounding-box filter (only exclude $[0,0,0,0]$)
  \item Corrected link label generation (one positive per GT KVP pair)
  \item $\lambda = 5.0$ (increased to compensate for sparser but correct supervision)
  \item Batch size 1 with 8 gradient-accumulation steps (effective batch size 8)
  \item Early-stopping patience 10 based on validation entity F1
\end{itemize}
Training targeted 30 epochs across resumed jobs; the final run history does not
support reporting a single uninterrupted epoch number for the selected
checkpoint.

\subsection*{Result 8: Regular-KVP F1 = 0.334 (text-only) / 0.309 (text+location)}
\label{sec:res8}

\begin{table}[ht]
  \centering
  \caption{Experiment~8: pair-level macro F1 under the released KVP10k
  benchmark. V4 emits regular linked pairs only; the All row additionally
  includes unkeyed and unvalued ground truth.}
  \label{tab:exp8-results}
  \begin{tabular}{@{}llccc@{}}
    \toprule
    \textbf{Category} & \textbf{System} & \textbf{Text} & \textbf{Location} & \textbf{Text+location} \\
    \midrule
    Regular & KVP10k paper & 0.659 & 0.650 & 0.611 \\
            & V4 (ours)    & 0.334 & 0.456 & 0.309 \\
    \midrule
    All     & KVP10k paper & 0.643 & 0.661 & 0.612 \\
            & V4 (ours)    & 0.214 & 0.310 & 0.199 \\
    \bottomrule
  \end{tabular}
\end{table}

With all three bugs corrected, V4 achieves regular-KVP F1 $=0.334$ for text,
$0.456$ for location, and $0.309$ for their combination. The corresponding
published scores are $0.659$, $0.650$, and $0.611$. V4 therefore remains
$0.325$ below the paper on text and $0.302$ below it on the combined metric.
The model emits only linked regular pairs, so it necessarily scores zero on
unkeyed and unvalued categories; including those categories lowers its All F1
to $0.214$ text-only and $0.199$ combined. A post-hoc decode-time recovery of
unlinked spans raises All F1 to $0.257$/$0.237$ without retraining and leaves
the regular scores unchanged (Analysis~C, Table~\ref{tab:alltype-recovery}).
These results are obtained with a
125M-parameter discriminative model trained on 55\% of the available training
data. They show that encoder-based linking is feasible, but do not support
parity with the published generative baseline. The entity classifier retains
F1 $=0.827$; two internal diagnostics (Section~\ref{sec:linker-diagnostics})
investigate why regular-pair linking remains weaker.

\paragraph{Note on the two metrics.}
Entity F1 and link F1 measure two different tasks and are not directly
comparable. \emph{Entity F1} measures whether the model correctly labels each
token as key, value, or other (finding the pieces), whereas \emph{link F1}
measures the harder task of pairing each detected key with its correct value
(connecting the pieces). A high entity F1 of 0.847 therefore does not imply a
high link F1: entity detection is comparatively strong, while linking remains
difficult. The full V4 model reports entity F1 $= 0.827$, slightly below the
0.847 of entity-only Stage~4a, because joint training with the linker trades a
small amount of entity accuracy for the linking objective.

\paragraph{Per-cluster breakdown.}
\label{sec:stage4-per-cluster}

\begin{table}[ht]
  \centering
  \caption{Experiment~8: V4 regular-KVP macro F1 by page-level
  annotation-geometry cluster (same benchmark matching as
  Table~\ref{tab:exp8-results}).}
  \label{tab:exp8-cluster}
  \begin{tabular}{@{}lrrrr@{}}
    \toprule
    \textbf{Stratum} & \textbf{Assigned} & \textbf{Scored} & \textbf{Text F1} & \textbf{Text+location F1} \\
    \midrule
    Cluster 0 (higher count, broad spread) & 319 & 304 & 0.290 & 0.264 \\
    Cluster 1 (lower count, compact)       & 213 & 179 & 0.406 & 0.383 \\
    Geometry unavailable                    &  49 &   0 & ---   & ---   \\
    \bottomrule
  \end{tabular}
\end{table}

V4 performs better on Cluster~1 under both regular-pair modes. This association
does not show that a semantic document type is easier: the clusters are based
only on annotation-box geometry, and no causal mechanism is tested. The 49
geometry-unavailable pages have no scored regular ground truth. Recomputing the
breakdown changes neither predictions nor the overall F1 values in
Table~\ref{tab:exp8-results}.

\subsection*{Effect of Matching Thresholds}
\label{sec:threshold-effect}

To separate text and location sensitivity, the same fixed V4 predictions are
re-scored under three benchmark-compatible variants. Table~\ref{tab:threshold-grid}
uses regular-KVP macro F1 throughout; no threshold is tuned on the test set.

\begin{table}[ht]
  \centering
  \caption{V4 regular-KVP macro F1 for identical predictions under diagnostic
  matching variants. The official row follows the released code's
  NED $<0.2$ and IoU $\geq0.3$ comparisons.}
  \label{tab:threshold-grid}
  \begin{tabular}{@{}lccc@{}}
    \toprule
    \textbf{Setting} & \textbf{NED} & \textbf{IoU} & \textbf{F1} \\
    \midrule
    Official combined & $<0.2$ & $\geq0.3$ & 0.309 \\
    Relaxed location  & $<0.2$ & $\geq0.1$ & 0.321 \\
    Relaxed text      & $<0.3$ & $\geq0.3$ & 0.321 \\
    Text only         & $<0.2$ & ---        & 0.334 \\
    \bottomrule
  \end{tabular}
\end{table}

Relaxing either axis yields almost the same gain ($+0.013$), while removing the
location condition entirely adds $0.025$. Boundary-near text and location
errors therefore contribute comparably, but neither explains most of the gap
to the published regular-KVP scores. The large residual consists of incorrect
or missing relations rather than predictions narrowly failing one acceptance
threshold.

% ---------------------------------------------------------------------------
\section{Linker Diagnostic Analyses}
\label{sec:linker-diagnostics}
% ---------------------------------------------------------------------------

To understand why the corrected V4 link F1 plateaus below the published
baseline, we run two post-hoc analyses on the trained V4 model. These are
\emph{diagnostics} on an existing model, not new trained variants.

\paragraph{Analysis A: Decision-threshold sensitivity.}
We sweep the link decision threshold on a held-out validation slice (a fixed
10\% of the training pages, never used for linker updates) from 0.05 to 0.90.
Link F1 is essentially flat across the entire range (text-only F1 stays near
0.34; text+location near 0.26), and the number of recovered true links does
not change materially. Recall is therefore not throttled by the score cut: it
is bounded by which candidate pairs the linker generates at all. Threshold
tuning is not a lever for this model.

\paragraph{Analysis B: Oracle entities (upper bound).}
We replace the predicted entity spans with the ground-truth key/value spans and
re-run linking, isolating linker quality from entity-detection quality.
Table~\ref{tab:oracle-linking} reports the legacy pooled regular-link diagnostic,
which is retained to compare end-to-end and oracle conditions within one
scorer but is not directly comparable with the macro benchmark tables. Perfect
entities raise F1 by only $\sim$0.02--0.05 in that diagnostic, indicating that
relation modeling remains the larger internal bottleneck.

\begin{table}[ht]
  \centering
  \caption{Analysis~B: end-to-end vs.\ oracle-entity linking under the legacy
  pooled regular-link diagnostic. These values are not KVP10k macro scores.}
  \label{tab:oracle-linking}
  \begin{tabular}{@{}lcc@{}}
    \toprule
    \textbf{Setting} & \textbf{End-to-end} & \textbf{Oracle entities} \\
    \midrule
    Link F1 (text-only)     & 0.342 & 0.365 \\
    Link F1 (text+location) & 0.309 & 0.345 \\
    \bottomrule
  \end{tabular}
\end{table}

Analyses~A and~B together indicate that the residual gap to the published
baseline is a \emph{relation-modeling} limitation rather than a threshold or
entity-detection artifact. This motivates the future-work direction of
segment-graph linkers such as SPADE and BROS
(Section~\ref{sec:limitations}).

\paragraph{Analysis C: Decode-time recovery of unkeyed and unvalued spans.}
The primary V4 model emits only linker-consumed regular pairs. The entity
classifier is nonetheless trained on every key and value span --- including
those belonging to unvalued keys and unkeyed values --- so these spans are
detected but discarded at decode time. As a post-hoc decoding change (no
retraining), we additionally emit each detected key span with no
above-threshold link as an \emph{unvalued} prediction and each value span never
selected by a regular pair as an \emph{unkeyed} prediction. The regular
predictions are byte-identical to the headline export (2{,}601 pairs), plus
1{,}211 recovered unvalued and 1{,}747 unkeyed predictions.
Table~\ref{tab:alltype-recovery} reports the resulting macro F1: regular scores
are unchanged by construction, while All F1 rises from $0.214/0.199$
(text / text+location) to $0.257/0.237$. The recovered categories are weak in
absolute terms --- unkeyed values are genuinely floating text with no anchoring
key --- so we retain the regular-only configuration as the headline and present
this recovery only as a lower-cost coverage extension.

\begin{table}[ht]
  \centering
  \caption{Analysis~C: macro F1 with decode-time recovery of unkeyed/unvalued
  spans (no retraining). Regular is the headline configuration and is unchanged
  by construction; ``All (regular only)'' is the primary model and ``All (with
  recovery)'' adds the recovered spans.}
  \label{tab:alltype-recovery}
  \begin{tabular}{@{}lccc@{}}
    \toprule
    \textbf{Category} & \textbf{Text} & \textbf{Location} & \textbf{Text+location} \\
    \midrule
    Regular (headline)     & 0.334 & 0.456 & 0.309 \\
    Unvalued (recovered)   & 0.271 & 0.369 & 0.242 \\
    Unkeyed (recovered)    & 0.178 & 0.325 & 0.165 \\
    All (regular only)     & 0.214 & 0.310 & 0.199 \\
    All (with recovery)    & 0.257 & 0.392 & 0.237 \\
    \bottomrule
  \end{tabular}
\end{table}

% ---------------------------------------------------------------------------
\section{Summary}
\label{sec:narrative-summary}
% ---------------------------------------------------------------------------

Table~\ref{tab:full-progression} summarises the full progression of
experiments.

\begin{table}[ht]
  \centering
  \caption{Full experimental progression. Published, Mistral, and V4 rows use
  regular-KVP macro F1 (text-only / text+location); V1 and V2 retain pooled
  regular-link diagnostic F1 and are not directly comparable.}
  \label{tab:full-progression}
  \small
  \begin{tabular}{@{}lccl@{}}
    \toprule
    \textbf{Model} & \textbf{Entity F1} & \textbf{Link F1} & \textbf{Key Finding} \\
    \midrule
    KVP10k paper baseline     & ---   & 0.659 / 0.611 & Published generative system \\
    Exp 1: Mistral-7B        & ---   & 0.598 / 0.521 & Re-implementation; no unkeyed outputs \\
    Exp 2: Stage~4a (entity only) & 0.847 & --- & Strong entity classification \\
    Exp 3: V1 (token linking) & 0.841 & 0.008 & 450:1 imbalance prevents linking \\
    Exp 4: V2 (span linking)  & 0.847 & 0.007 & Bugs mask span-level benefit \\
    Exp 5--7: Diagnostics     & ---   & ---   & 3 pipeline bugs identified \& fixed \\
    Exp 8: V4 (corrected)     & 0.827 & 0.334 / 0.309 & Regular KVPs only; below baseline \\
    \bottomrule
  \end{tabular}
\end{table}

This section summarises the experimental progression. The work began with the
question of whether a layout-aware discriminative model could match a
generative reference baseline that had not been benchmarked against encoder
models. The generative approach is competitive overall but can produce
unsupported text on some documents, whereas token-grounded discriminative
models cannot generate text absent from the OCR input. Making the
discriminative model link key-value pairs proved non-trivial: token-level
linking is hindered by severe class imbalance, span-level linking exposed
pipeline failures, and non-trivial V4 results emerged after all three
corrections were applied jointly. The final 125M-parameter model reaches
regular-KVP macro F1 $=0.334$ for text and $0.309$ for text+location, below the
paper's $0.659$ and $0.611$. The independent numerical effect of each
correction was not isolated in this study. Internal diagnostics localise the
remaining gap primarily to relation modeling.
